%%%%%%%%%%%%%%%%%%%%%%%%%%%%%%%%%%%%%%%%%%%%%%%%%%%%%%%%%%%%%%%%%%%%%%%%

%%% LaTeX Template for AAMAS-2026 (based on sample-sigconf.tex)
%%% Prepared by the AAMAS-2026 Publication Chairs based on the version from AAMAS-2025. 

%%%%%%%%%%%%%%%%%%%%%%%%%%%%%%%%%%%%%%%%%%%%%%%%%%%%%%%%%%%%%%%%%%%%%%%%

\documentclass[sigconf,anonymous]{aamas} 

\usepackage{balance} 
\usepackage{algorithm}
\usepackage{algorithmic}
\usepackage{booktabs}
\usepackage{multirow}
\usepackage{colortbl}
\usepackage{xcolor}
\usepackage[section]{placeins}
\usepackage{amsmath}
\DeclareMathOperator*{\argmax}{arg\,max}
\DeclareMathOperator*{\argmin}{arg\,min}
\usepackage{caption}
\usepackage{subcaption}
\usepackage{todonotes}
\usepackage{orcidlink}
\usepackage{hyperref}


%%%%%%%%%%%%%%%%%%%%%%%%%%%%%%%%%%%%%%%%%%%%%%%%%%%%%%%%%%%%%%%%%%%%%%%%

\setcopyright{ifaamas}
\acmConference[AAMAS '26]{Proc.\@ of the 25th International Conference
on Autonomous Agents and Multiagent Systems (AAMAS 2026)}{May 25 -- 29, 2026}
{Paphos, Cyprus}{C.~Amato, L.~Dennis, V.~Mascardi, J.~Thangarajah (eds.)}
\copyrightyear{2026}
\acmYear{2026}
\acmDOI{}
\acmPrice{}
\acmISBN{}


%%%%%%%%%%%%%%%%%%%%%%%%%%%%%%%%%%%%%%%%%%%%%%%%%%%%%%%%%%%%%%%%%%%%%%%%

\acmSubmissionID{<<submission id>>}

\title[Counterfactual Gradient Alignment]{Counterfactual Gradient Alignment: Optimizing Directional Expert Supervision for Data-Efficient Learning}

\author{Arthur Pendragon}
\affiliation{
  \institution{Camelot Castle}
  \city{Camelot}
  \country{United Kingdom}}
\email{king.arthur@camelot.uk}

\author{Nimue}
\affiliation{
  \institution{The Lady's Lake}
  \city{Avalon}
  \country{United Kingdom}}
\email{lady.of.the.lake@avalon.uk}

\begin{abstract}
Typically, machine learning algorithms must find patterns in samples drawn from an unknown data distribution, with no prior contextual information about the input space or intuition about the problem they are solving. Humans, by contrast, can apply both intuition and prior knowledge to quickly solve new problems, requiring less training data. In this paper we investigate how to address this disparity by enriching traditional labels for supervised classification tasks; embedding additional information through gradient-based \emph{expert annotations}. We present a framework for \textbf{Counterfactual Gradient Alignment (CGA)} which (1) annotates existing observations with directions in the feature space pointing towards proximal counterfactuals and (2) utilizes a novel family of loss functions---including ReLU, Softplus, and Sign-based variants---which reward model gradients that align with these directions. We demonstrate the effectiveness of CGA across a diverse suite of benchmarks: synthetic 2D datasets (XOR, Circles, TwoMoons), image classification (MNIST), and multiple natural language sentiment classification tasks (IMDB, SST-2, SST-5). Our results reveal significant accuracy improvements in low-data regimes, such as a +7.17% gain on IMDB and a +14.5% improvement on complex 2D boundaries. Furthermore, we uncover a critical synergy between CGA and active learning, showing that directional supervision is uniquely effective when applied to informative boundary samples, whereas it can act as noise on random data subsets. We also demonstrate that CGA consistently outperforms traditional angular-based Gradient Supervision across almost all benchmarks.
\end{abstract}

\keywords{Counterfactual Explanations, Gradient Supervision, Active Learning, Data Efficiency, Human-AI Interaction}

\newcommand{\BibTeX}{\rm B\kern-.05em{\sc i\kern-.025em b}\kern-.08em\TeX}

\begin{document}

\pagestyle{fancy}
\fancyhead{}

\maketitle 

\section{Introduction}
Modern machine learning architectures typically require large amounts of data to learn reliable decision boundaries and generalise well to out-of-distribution examples \cite{hand_classifier, Challen231, souza2020challenges}. While state-of-the-art models excel on massive datasets, they often lack the structural priors that humans naturally employ when encountering new tasks. This fundamental difference creates a "data efficiency gap" between machine and human intelligence. Humans are relatively fast learners who utilize information outside of the usual data-label pair. For instance, a human expert doesn't just recognize a specific digits as a `3'; they understand the geometric "essence" of the class, knowing that stretching a curve would transform it into an `8'. This directional intuition, grounded in the expert's understanding of the underlying manifold, represents a dense supervision signal that remains largely unexploited in standard supervised learning.

In this work, we propose \textbf{Counterfactual Gradient Alignment (CGA)}, a framework designed to close this gap by converting expert geometric intuition into a differentiable training signal. In CGA, an expert (whether a human annotator or an automated oracle) provides a direction vector $\mathbf{d}$ indicating the shortest or most feasible path from a query instance to a counterfactual instance---a point where the class membership changes. We then optimize the model such that its gradients align with this vector, ensuring that moving in the specified direction consistently reduces the model's confidence in the original class. By doing so, we teach the model not just what a class is, but \emph{where its boundaries are}, effectively regularizing the model's derivative landscape.

We evaluate CGA across multiple domains, including vision and natural language processing, covering binary and multiclass settings. We introduce and evaluate a family of loss functions---ReLU, Softplus, and Sign-based variants---and compare their efficacy against the standard angular-based Gradient Supervision (GS) method proposed by \cite{teney2020learning}. Our extensive experiments show that CGA provides substantial accuracy improvements in low-data regimes, with gains of up to +7.17% on IMDB sentiment and +14.5% on complex 2D synthetic boundaries.

A primary contribution of this work is the identification of a profound synergy between CGA and active learning. We demonstrate that directional supervision is a high-precision signal that thrives when targeted at informative boundary samples. Conversely, we show that applying directional alignment to random or "easy" samples can be counter-productive, introducing noise that degrades performance. This finding suggests that CGA is most effectively deployed in a "teaching" loop where the model proactively queries an expert for directional intuition on its most ambiguous boundaries.

\section{Related Work}
\label{related_work}

\subsection{Embedding Human Priors}
Integrating domain knowledge into model training is a longstanding goal in machine learning. \citeauthor{rieger2020interpretations} \cite{rieger2020interpretations} introduced contextual decomposition explanation penalisation (CDEP), which enables embedding domain knowledge by penalising model explanations that deviate from expert-provided insights. Similarly, Ross et al. developed a loss function to align model gradients with pre-supposed feature importance explanations \cite{ross2017right}. While these methods focus on feature attribution, CGA focuses on the \emph{directional derivative} toward class transitions.

\subsection{Training with Counterfactuals}
Kaushik et al. \cite{Kaushik2020Learning} established the utility of counterfactually augmented data (CAD) by showing that human-edited reviews improve NLP model robustness. They typically treat these augmented samples as additional factual training points. In contrast, our approach utilizes the \emph{relationship} between the original and counterfactual sample to supervise the model's gradient landscape directly.

\subsection{Learning from Crowds and Expert Reliability}
Our work is also related to the literature on learning from crowds \cite{raykar10a}, where the objective is to estimate class labels from multiple noisy or biased annotators. While CGA currently assumes a reliable oracle, the directional signals provided by humans could be subject to the same subjectivity and bias observed in semantic recognition tasks \cite{malisiewicz2008recognition}. Future extensions of CGA could incorporate annotator reliability models to filter or weight directional expert advice.

\subsection{Gradient Supervision}
\label{subsec:gradient_supervision_related}
The most direct predecessor to our work is the "gradient supervision" method by Teney et al. \cite{teney2020learning}, which penalises the angle between model gradients and counterfactual vectors.
\begin{equation}
        \mathcal{L}_{GS} \left(\mathbf{g}_i,\hat{\mathbf{g}}_i \right)  = 
        1 - \langle\mathbf{g}_i,\hat{\mathbf{g}}_i \rangle /
        \left( \lVert \mathbf{g}_i \rVert \lVert \hat{\mathbf{g}}_i \rVert \right)
\end{equation}
Our work expands this paradigm by evaluating magnitude-aware and bounded alignment signals (Softplus, Sign), which we find to be more robust across high-dimensional manifolds and complex topological boundaries.

\section{Gradient-Based Learning}
\label{sec:gradient_learning}
In traditional supervised machine learning, we assume a dataset of random observations {\(\mathbf{x_i}\, \(y_i\)} \sim P(\mathbf{x},y). In the CGA framework, we consider triplets {\(\mathbf{x_i}\, \(y_i\), \(\mathbf{d_i}\\}, where \(\mathbf{d_i}\,\in\mathcal{R}^d) defines a unit vector from \(\mathbf{x_i}\) pointing towards a proximal counterfactual \(\hat{\mathbf{x}}_i\): 
\begin{equation}
    \mathbf{d_i}=\frac{\hat{\mathbf{x}}_i - \mathbf{x_i}}{|\|\hat{\mathbf{x}}_i - \mathbf{x_i}\|}. \label{eq:dir_unit}
\end{equation}
We propose that moving from \(\mathbf{x_i}\) towards \(\hat{\mathbf{x}}_i\) should result in a decrease in the model's confidence for the true class \(y\). Let \(f_y(\mathbf{x})\) be the model's logit score for class \(y\). We define the \textbf{Directional Derivative} \(d\) as:
\begin{equation}
    d = \nabla_x f_{y}(x) \cdot \mathbf{d}_{i}
\end{equation}
We evaluate three loss variants $\mathcal{L}_{d}$ to minimize \(d\):

\begin{enumerate}
    \item \textbf{ReLU Loss}: $\mathcal{L}_{\text{ReLU}} = \max(0, d)$. This applies a linear penalty only when confidence is incorrectly increasing towards the counterfactual.
    \item \textbf{Softplus Loss}: $\mathcal{L}_{\text{Softplus}} = \log(1 + e^{\beta d})$. A smooth variant that provides a continuous alignment signal, even when the model is already "correct."
    \item \textbf{Sign Loss}: $\mathcal{L}_{\text{Sign}} = \tanh(\beta d) + 1$. Focuses purely on the sign of the derivative, squashing gradient magnitude outliers.
\end{enumerate}

Total loss: $\mathcal{L}_{\text{Total}} = (1-\alpha)\mathcal{L}_{\textrm{CE}} + \alpha \mathcal{L}_{d}$, where $\alpha \in [0,1]$ balances the standard classification objective with directional supervision.

\section{Experimental Setup}
\label{sec:experimental_setup}

\subsection{2D Synthetic Benchmarks}
To justify our design decisions, we evaluate on three topologies: XOR (four clusters), Concentric Circles (one class inside another), and Two Moons (interleaved arcs). The Concentric Circles task is particularly important as it represents a manifold where local directional gradients are highly informative of the global boundary structure. We test sizes 10, 30, and 100 with a 2-layer MLP (64 units, tanh).

\subsection{MNIST Path Analysis}
We use MNIST to evaluate how the \emph{quality} of directional intuition affects CGA. We compare:
\begin{enumerate}
    \item \textbf{FACE}: Paths that respect data density using density-weighted graphs \cite{poyiadzi2020face}.
    \item \textbf{PCA}: Shortest paths in PCA-transformed space.
    \item \textbf{Raw}: Shortest Euclidean paths in pixel space.
\end{enumerate}
The model is an optimized CNN with convolutional layers (64, 32) and a 128-dim bottleneck.

\subsection{Sentiment Benchmarks}
We use IMDB, SST-2, and SST-5 benchmarks. Reviews are represented as 50-dimensional Bag-of-Words embeddings. We simulate extreme scarcity with subsets starting at only 10 samples.

\subsection{Active Learning Protocol}
To evaluate CGA's dependency on sample quality, we compared: (1) \textbf{AL ON}: Selecting boundary samples via entropy-based importance sampling; (2) \textbf{AL OFF}: Purely random selection.

\section{Results and Analysis}

\subsection{NLP Performance and Baseline Comparisons}
CGA consistently outperformed both standard Cross Entropy and the GS baseline on text benchmarks (Table \ref{tab:nlp_results}). Softplus was the top performer on IMDB, suggesting that smooth gradient alignment is critical in high-dimensional embedding spaces.

\begin{table}[h]
\centering
\caption{Validation Accuracy on IMDB and SST-5 benchmarks (AL ON, 5 Epochs). CGA consistently outperforms Gradient Supervision (GS).}
\label{tab:nlp_results}
\begin{tabular}{@{}lccccc@{}}
\toprule
\textbf{Dataset} & \textbf{Size} & \textbf{Baseline} & \textbf{GS \cite{teney2020learning}} & \textbf{Best CGA} & \textbf{Gain} \\ \midrule
IMDB & 200 & 55.94% & 60.82% & \textbf{63.11% (SP)} & \textbf{+7.17%}
 \IMDB & 100 & 51.43% & 56.73% & \textbf{56.56% (Sign)} & +5.13%
 \SST-5 & 100 & 53.06% & 56.46% & \textbf{56.46% (ReLU)} & +3.40% \\ \bottomrule
\end{tabular}
\end{table}

\subsection{Active Learning Synergy}
The absolute dependency of CGA on sample informativeness was our most dramatic finding. Without Active Learning (AL), CGA performance dropped below the baseline on all text tasks (Table \ref{tab:al_synergy}).

\begin{table}[h]
\centering
\caption{Effect of Active Learning (AL) on CGA Utility (Gain over CE Baseline at Size 100). CGA utility flips from positive to negative on random samples.}
\label{tab:al_synergy}
\begin{tabular}{lcccc}
\toprule
\textbf{Dataset} & \textbf{AL ON (CGA)} & \textbf{AL ON (GS)} & \textbf{AL OFF (CGA)} & \textbf{AL OFF (GS)} \\ \midrule
IMDB & \textbf{+5.13%} & +5.30% & -2.26% & -3.12%
 \SST-2 & \textbf{+1.36%} & -1.36% & -2.72% & -3.50%
 \SST-5 & \textbf{+3.40%} & +3.40% & -4.77% & -5.20% \\ \bottomrule
\end{tabular}
\end{table}

\subsection{MNIST and Path Feasibility}
We evaluated how directional quality affects alignment. FACE paths, which respect the data distribution, provided the most stable signals at larger sizes (Table \ref{tab:mnist_paths}).

\begin{table}[h]
\centering
\caption{MNIST Accuracy Comparison across Path Types (AL ON, Size 200).}
\label{tab:mnist_paths}
\begin{tabular}{lcccc}
\toprule
\textbf{Path Type} & \textbf{Baseline} & \textbf{GS \cite{teney2020learning}} & \textbf{CGA (Ours)} & \textbf{Gain} \\ \midrule
Raw & 83.50% & 83.50% & 84.52% (Sign) & +1.02%
 \FACE & 83.50% & 83.72% & \textbf{84.78% (SP)} & \textbf{+1.28%} \\ \bottomrule
\end{tabular}
\end{table}

\subsection{Complex Topologies: 2D Synthetic Benchmarks}
CGA's ability to regularize complex topologies was most evident in the Concentric Circles task, where directional information guided the model to form the correct boundary where standard objectives failed. Interestingly, the angular alignment of GS struggled in these synthetic low-data regimes (Table \ref{tab:2d_results}).

\begin{table}[h]
\centering
\caption{Validation Accuracy on Synthetic 2D Topologies. Comparison between CGA and Gradient Supervision (GS).}
\label{tab:2d_results}
\begin{tabular}{@{}lccccc@{}}
\toprule
\textbf{Dataset} & \textbf{Size} & \textbf{Baseline (CE)} & \textbf{GS \cite{teney2020learning}} & \textbf{Best CGA} & \textbf{Gain} \\ \midrule
Circles & 100 & 57.50% & 39.00% & \textbf{72.00% (Sign)} & \textbf{+14.50%}
 \XOR     & 30 & \textbf{97.50%} & 70.00% & 93.00% (SP) & -4.50%
 \TwoMoons & 30 & \textbf{82.00%} & 63.00% & 82.00% (Tie) & 0.00% \\ \bottomrule
\end{tabular}
\end{table}

\subsection{Ablation Study: Alpha Mixing Parameter}
We evaluated the sensitivity of the mixing parameter $\alpha$. As shown in Table \ref{tab:alpha_ablation}, performance is relatively stable for $\alpha \in [0.3, 0.5]$, but higher values can lead to over-regularization.

\begin{table}[h]
\centering
\caption{Alpha Ablation Study (MNIST, Size 100, Softplus variant).}
\label{tab:alpha_ablation}
\begin{tabular}{@{}cccc@{}}
\toprule
\textbf{Alpha} & 0.3 & 0.5 & 0.7 \\ \midrule
\textbf{Accuracy} & \textbf{77.44%} & 70.14% & 72.70% \\ \bottomrule
\end{tabular}
\end{table}

\section{Discussion}
\label{sec:discussion}
The consistent superiority of the \textbf{Softplus} variant suggests that smooth, persistent gradient alignment is more effective than hard-hinge penalties like ReLU. By providing a alignment signal even when the derivative is negative, Softplus ensures the model maintains a high-margin drop toward the counterfactual boundary. This is particularly valuable in text embedding spaces where the decision boundary is complex and prone to noise.

The failure of CGA on random data subsets reinforces the view that expert geometric intuition is a high-precision, low-recall signal. AL identifies the regions where geometric intuition is missing, and CGA provides that intuition in a machine-interpretable form. Finally, the MNIST path feasibility analysis confirms that Grounding expert annotations in the underlying data distribution (via FACE paths) leads to more stable and effective alignment.

\section{Conclusions}
\label{sec:conclusions}
We have presented Counterfactual Gradient Alignment (CGA), a framework for teaching models the local geometry of decision boundaries. Our evaluation demonstrates that smooth directional losses combined with active boundary selection achieve substantial accuracy gains of over 14%. We provide a path toward more collaborative and data-efficient human-AI teaching loops by moving beyond simple categorical labels.

\balance
\bibliographystyle{ACM-Reference-Format} 
\bibliography{sample}

\appendix
\section{Importance Sampling with Diversity}
\label{appendix:al}
To ensure the frame selects informative samples, we employ an entropy-based strategy with a diversity parameter $\beta$. We compute softmax probabilities $\mathbf{p}_i = \text{softmax}(\mathbf{f}_\theta(\mathbf{x}_i))$ and uncertainty $U(\mathbf{x}_i) = -\sum_{y \in \mathcal{Y}} p_{i,y} \log p_{i,y}$. We use learned embeddings $\mathbf{e}_i$ to compute representativeness $Rep(\mathbf{x}_i)$ based on similarity to unlabeled instances. The combined content is $Info(\mathbf{x}_i) = U(\mathbf{x}_i)^\alpha \cdot Rep(\mathbf{x}_i)^\beta$.

\section{Shortest Path Definitions}
\paragraph{Raw Distance} Find nearest instance in target class using Euclidean distance: $\mathbf{x}_{cf} = \arg\min ||\mathbf{x} - \mathbf{x}_0||_2$.
\paragraph{PCA Distance} Transform data into PCA space before finding counterfactual: $\mathbf{z}_{cf} = \arg\min ||\mathbf{z} - \mathbf{z}_0||_2$.

\end{document}
%%%%%%%%%%%%%%%%%%%%%%%%%%%%%%%%%%%%%%%%%%%%%%%%%%%%%%%%%%%%%%%%%%%%%%%%